\documentclass[a4paper]{article}
\usepackage{geometry}
\usepackage{graphicx}
\usepackage{natbib}
\usepackage{amsmath}
\usepackage{amssymb}
\usepackage{amsthm}
\usepackage{paralist}
\usepackage{epstopdf}
\usepackage{tabularx}
\usepackage{longtable}
\usepackage{multirow}
\usepackage{multicol}
\usepackage[hidelinks]{hyperref}
\usepackage{fancyvrb}
\usepackage{algorithm}
\usepackage{algorithmic}
\usepackage{float}
\usepackage{paralist}
\usepackage[svgname]{xcolor}
\usepackage{enumerate}
\usepackage{array}
\usepackage{times}
\usepackage{url}
\usepackage{fancyhdr}
\usepackage{comment}
\usepackage{environ}
\usepackage{times}
\usepackage{textcomp}
\usepackage{caption}
\usepackage{color}
\usepackage{xcolor}



\urlstyle{rm}

\setlength\parindent{0pt} % Removes all indentation from paragraphs
\theoremstyle{definition}
\newtheorem{definition}{Definition}[]
\newtheorem{conjecture}{Conjecture}[]
\newtheorem{example}{Example}[]
\newtheorem{theorem}{Theorem}[]
\newtheorem{lemma}{Lemma}
\newtheorem{proposition}{Proposition}
\newtheorem{corollary}{Corollary}

\floatname{algorithm}{Procedure}
\renewcommand{\algorithmicrequire}{\textbf{Input:}}
\renewcommand{\algorithmicensure}{\textbf{Output:}}
\newcommand{\abs}[1]{\lvert#1\rvert}
\newcommand{\norm}[1]{\lVert#1\rVert}
\newcommand{\RR}{\mathbb{R}}
\newcommand{\CC}{\mathbb{C}}
\newcommand{\Nat}{\mathbb{N}}
\newcommand{\br}[1]{\{#1\}}
\DeclareMathOperator*{\argmin}{arg\,min}
\DeclareMathOperator*{\argmax}{arg\,max}
\renewcommand{\qedsymbol}{$\blacksquare$}

\definecolor{dkgreen}{rgb}{0,0.6,0}
\definecolor{gray}{rgb}{0.5,0.5,0.5}
\definecolor{mauve}{rgb}{0.58,0,0.82}

\newcommand{\Var}{\mathrm{Var}}
\newcommand{\Cov}{\mathrm{Cov}}

\newcommand{\vc}[1]{\boldsymbol{#1}}
\newcommand{\xv}{\vc{x}}
\newcommand{\Sigmav}{\vc{\Sigma}}
\newcommand{\alphav}{\vc{\alpha}}
\newcommand{\muv}{\vc{\mu}}

\newcommand{\red}[1]{\textcolor{red}{#1}}

\def\x{\mathbf x}
\def\y{\mathbf y}
\def\w{\mathbf w}
\def\v{\mathbf v}
\def\E{\mathbb E}
\def\V{\mathbb V}

% TO SHOW SOLUTIONS, include following (else comment out):
\newenvironment{soln}{
    \leavevmode\color{blue}\ignorespaces
}{}

\newcommand \mle [1]{{\hat #1}^{\rm MLE}}



\hypersetup{
%    colorlinks,
    linkcolor={red!50!black},
    citecolor={blue!50!black},
    urlcolor={blue!80!black}
}

\geometry{
  top=1in,            % <-- you want to adjust this
  inner=1in,
  outer=1in,
  bottom=1in,
  headheight=3em,       % <-- and this
  headsep=2em,          % <-- and this
  footskip=3em,
}


\pagestyle{fancyplain}
\lhead{\fancyplain{}{Homework 4}}
\rhead{\fancyplain{}{CS 760 Machine Learning}}
\cfoot{\thepage}

\newcommand{\R}{\mathbb R}
\newcommand{\sgn}{\mathrm{sign}}
\usepackage{enumitem}
\usepackage{bm}
\def\*#1{\mathbf{#1}}

\title{\textsc{Homework 4}} % Title

%%% NOTE:  Replace 'NAME HERE' etc., and delete any "\red{}" wrappers (so it won't show up as red)

\author{
\red{$>>$NAME HERE$<<$} \\
\red{$>>$ID HERE$<<$}\\
} 

\date{}

\begin{document}

\maketitle 


\textbf{Instructions:} 

Use this LaTeX file as a template to develop your homework and submit it as a single PDF file on Gradescope. 
Remember to assign the pages to problems when submitting.

\section{Verify PAC Bound}
Let $h$ be the VC dimension of function family $F$.
For any $\delta>0$, with probability at least $1-\delta$ we have
$$R(\hat f_S) - \hat R_S(\hat f_S) \le 2 \sqrt{\frac2n\left(h \log n + h \log\frac{2e}{h}  + \log \frac{2}{\delta}\right)}$$
where the log is natural, $R(f) = \E 1_{[f(x)\neq y]}$ is the risk of $f$,
$S=(x_1,y_1), \ldots, (x_n, y_n)$ is a training set of size $n$,
$\hat R_S(f) = {1\over n} \sum_{i=1}^n 1_{[f(x_i)\neq y_i]}$ is the empirical risk of $f$ on $S$,
and 
$\hat f_S \in \argmin_{f\in F} \hat R_S(f)$ is an empirical risk minimizer~(ERM).
We now verify the PAC bound on a simple classification task.  Let $p(x)=\mathrm{uniform}([-1,1])$.  Given $x$, the label is deterministic: $y=\sgn(x)$.
Recall $\sgn(x)=1$ if $x\ge 0$, and 0 otherwise.
This means the true decision boundary is at $x=0$.
Let $f_\theta(x) := \sgn(x-\theta)$ which has threshold at $\theta$.
Let $F=\{f_\theta: \theta \in [-1,1]\}$.  

\subsection{ERM [5 points]}

Given $S$, suppose we find the smallest positive item: $a = \min_{(x_i, y_i)\in S: y_i=1} x_i$ if one exists, otherwise let $a=1$.  Is $\hat f_S := f_a$ an ERM?  Justify your answer.

{\color{blue}
	Your solution here.
}

\subsection{True risk [5 points]}
Derive $R(f_a)$.

{\color{blue}
	Your solution here.
}

\subsection{Empirical risk [5 points]}
Derive $\hat R_S(f_a)$.

{\color{blue}
	Your solution here.
}

\subsection{VC dimension [10 points]}

What is the VC dimension of $F$?  Note $F$ contains threshold classifiers of the type ``left negative, right positive.''

{\color{blue}
	Your solution here.
}

\subsection{Generalization error [5 points]}

Fix $\delta=0.05$ (95\% confidence) and $n=200$.  What is the bound on the true risk of $\hat f_S$ given by the PAC bound?

{\color{blue}
	Your solution here.
}


\subsection{Empirical verification [10 points]}

Generate 10,000 random training sets from $p(x)$ and the associated labels.  Each training set $S$ has $n=200$ points.
On each $S$ you will compute $R(f_a) - \hat R_S(f_a)$.
Now you have 10,000 numbers: (1) produce a histogram of them, (2) find their 95\% quantile of them, and (3) compare the 95\% quantile to the bound in the previous question.  
Discuss your observations.

{\color{blue}
	Your solution here.
}

\section{VC dimension}

\subsection{Showing VC dimension [10 points]}
Let the input $x\in \mathcal{X}=\R$.
Consider $\mathcal{F}=\{f(x)=\sgn(ax^2+bx+c): a, b, c \in \R\}$, where $\sgn(z)=1$ if $z\ge0$, and 0 otherwise.
Compute the VC-dimension $VC(\mathcal{F})$ of $\mathcal F$.\\
{\color{blue}
	Your solution here.
}

\subsection{Shatter Coefficients [10 points]}
	The main intuition behind VC theory is that, although a collection of classifiers may be infinite, using a finite set of training data to select a good rule effectively reduces the number of different classifiers we need to consider. We can measure the effective size of class $\mathcal{F}$ using \textit{shatter coefficient}. Suppose we have a training set $D_n = \{(x_i, y_i)\}_{i=1}^n$ for a binary classification problem with labels $y_i = \{-1, +1\}$. Each classifier in $\mathcal{F}$ produces a binary label sequence
	$$\big(f(x_1), \cdots, f(x_n)\big) \in \{-1, +1\}^n$$
	There are at most $2^n$ distinct sequences, but often no all sequences can be generated by functions in $\mathcal{F}$. Shatter coefficient $\mathcal{S}(\mathcal{F}, n)$ is defined as the maximum number of labeling sequences the class $\mathcal{F}$ induces over $n$ training points in the feature space $\mathcal{X}$. More formally, 
	$$\mathcal{S}(\mathcal{F}, n) = \max_{x_1, \cdots, x_n \in \mathcal{X}}
	\Bigg|\Big\{\big(f(x_1), \cdots, f(x_n)\big)\Big\} \in \{-1, +1\}^n, f \in \mathcal{F}\Bigg|$$
	where $|\cdot|$ denotes the number of elements in the set.
	The \textit{Vapnik-Chervonenkis (VC) dimension} $V(\mathcal{F})$ of a class $\mathcal{F}$  is defined as the largest integer $k$ such that $\mathcal{S}(\mathcal{F}, k) = 2^k$. 
	
	Suppose there are $n$ points $(x_1, \cdots, x_n)\in\R$. Let $\mathcal{F}$ be the collection of 1-d linear classifiers: for each $t\in [0, 1], f_t \in \mathcal{F}$ labels $x\leq t$ as $-1$ and $x > t$ as $+1$, or vice-versa. What is the shatter coefficient $S(\mathcal{F}, n)$? What is VC-dimension $V(\mathcal{F})$? How can you get it from shatter coefficient?\\
{\color{blue}
	Your solution here.
}

\subsection{Monotonicity of VC-dimension [10 points]}
	Use the shatter coefficients to show the following monotonicity property of VC-dimension: if $\mathcal{F}$ and $\mathcal{F}'$ are hypothesis classes with $\mathcal{F} \subset \mathcal{F}'$, then the VC-dimension of $\mathcal{F}'$ is greater than or equal to that of $\mathcal{F}$.\\
	{\color{blue}
		Your solution here.
}


\section{Kernel SVM}
Consider the following kernel function defined over $z,z'\in Z$: $k(z,z') =
\begin{cases}
1 & \text{~if~} z=z', \\
0 & \text{~otherwise.}
\end{cases}$

\subsection{PSD [10 points]}

	Prove that for any integer $m>0$, any $z_1, \ldots, z_m \in Z$, the $m \times m$ kernel matrix $K=[K_{ij}]$ is positive semi-definite, where $K_{ij}=k(z_i, z_j)$ for $i,j=1\ldots m$. (Let us assume that for $i \neq j$, we have $z_i \neq z_j$.)
	Hint: An $m\times m$ matrix $K$ is positive semi-definite if $\forall v \in \R^m, v^\top K v \ge 0$. \\
	{\color{blue}
		Your solution here.
}
	
\subsection{Separability [10 points]}
	
	Given a training set $(z_1, y_1), \ldots, (z_n, y_n)$ with binary labels, the dual SVM problem with the above kernel $k$ will have parameters $\alpha_1, \ldots, \alpha_n, b \in \R$.  (Assume that for $i \neq j$, we have $z_i \neq z_j$.) The predictor for input $z$ is
	$$f(z) = \sum_{i=1}^n \alpha_i y_i k(z_i, z) + b.$$
	Recall the prediction is $\sgn(f(z))$.
	Prove that there exists $\alpha_1, \ldots, \alpha_n, b$ s.t. $f$ correctly separates the training set.\looseness-1\\
	{\color{blue}
		Your solution here.
}
	
\subsection{Prediction [10 points]}
	
	How does that $f$ predict input $z$ that is not in the training set? \\
	{\color{blue}
		Your solution here.
}

\end{document}
